(1) 用归一Jones态 \(|\psi(\theta,\phi)\rangle=(\cos(\theta/2),e^{i\phi}\sin(\theta/2))^{\mathsf T}\) 表示Poincaré球上的纯偏振。计算Berry联络 \(A_\phi=i\langle\psi|\partial_\phi\psi\rangle\)、\(A_\theta\) 和曲率 \(F_{\theta\phi}\)。
(2) 光的偏振态绝热沿固定 \(\theta\) 纬线绕行 \(\phi:0\to2\pi\)。求几何相位，并与该纬线所围球面立体角比较；特别计算赤道闭路。
(3) 两束等强相干光经历不同偏振演化但在输出具有相同Jones态，路径合成闭路围出有向立体角 \(\Omega\)。求探测强度随可调参考相位 \(\delta\) 的关系，说明如何区分几何相位与共同动力学相位。
(4) 对规范变换 \(|\psi\rangle\mapsto e^{i\chi(\theta,\phi)}|\psi\rangle\) 推导联络变换律，证明闭路可观测相位模 \(2\pi\) 不变，并解释单一规范在整球面不可光滑覆盖的原因。
